\reading{LTR-6}{Intraday Liquidity Risk Management}{DEFER}{0}

\begin{lrkeybox}[Learning objectives — official 2026 spine wording]
\begin{enumerate}[label=\textbf{\alph*.}]
\item Identify and explain the uses and sources of intraday liquidity.
\item Discuss the governance structure of intraday liquidity risk management.
\item Differentiate between methods for tracking intraday flows and monitoring risk levels.
\end{enumerate}
{\footnotesize\color{FRMGrey}Source: Shyam Venkat and Stephen Baird, \textit{Liquidity Risk Management: A Practitioner's Perspective}, Chapter 4.}
\end{lrkeybox}

% =====================================================================
\lo{LTR-6 a}{Identify and explain the uses and sources of intraday liquidity.}

\subsection{Uses of intraday liquidity}

\begin{itemize}
\item \term{Outgoing wire transfers}. Outgoing wire payments, on behalf of clients or the bank's own account, are typically the \textbf{largest use} of intraday liquidity.
\item \term{Settlements at Payment, Clearing and Settlement (PCS) systems}. Most PCS systems have one settlement per day, with many occurring in the late-afternoon timeframe.
\item \term{Funding of nostro accounts}. Cash transfer to a correspondent bank for services provided.
\item \term{Collateral pledging}.
\item \term{Asset purchases and funding}.
\end{itemize}

\subsection{Sources of intraday liquidity}

\begin{itemize}
\item \term{Cash balances}. Deposits at the central bank and at correspondent bank nostro accounts.
\item \term{Incoming funds flow}. Incoming flows from payments and \term{financial market utility (FMU)} settlements are the \textbf{largest source} of intraday funding in periods of normal market function.
\item \term{Intraday credit}. A credit line or overdraft permitted during business hours and covered by close of business. Lines are often \textbf{uncommitted} and provided \textbf{without interest charges}. Central banks serve as a large source of intraday credit for the banking system.
\item \term{Liquid assets}.
\item \term{Overnight borrowings}. Fed funds, Eurodollar borrowing, overnight deposits.
\item \term{Other term funding}. The bank would only do this if it has the appetite for longer-term funding (one week, one month, and so on).
\end{itemize}

\begin{lrtrapbox}[The two ``largest'' claims are a matched pair]
\textbf{Largest use} $=$ outgoing wire payments. \textbf{Largest source} $=$ incoming funds flow, and only \emph{in periods of normal market function}. Both superlatives are offered against each other, and the qualifier on the second one is droppable — a distractor stating that incoming flows are the largest source unconditionally has lost the scope condition that matters most, since the whole point of intraday liquidity risk is what happens when market function is \emph{not} normal.
\end{lrtrapbox}

\begin{lrtrapbox}[Intraday credit is uncommitted and free — until it is not]
Intraday credit lines are typically \textbf{uncommitted} (so they can be withdrawn) and carry \textbf{no interest charge} (so their cost does not appear in funding cost measures). A statement calling them committed, or interest-bearing, is a single-word corruption of a true claim.
\end{lrtrapbox}

\subsection{The treasurer's challenge}

The bank treasurer's challenge, managing the cacophony generated by high volumes of incoming and outgoing cash transactions, is affected by the following factors:

\begin{itemize}
\item \textbf{The variability of cash flow patterns.}
\item \textbf{The impact of market forces.}
  \begin{itemize}
  \item Daily volatility in asset prices can result in unanticipated margin calls that require additional cash funding.
  \item Central bank money desks that implement monetary policy directly may influence available liquidity in the marketplace, affecting a bank's ability to source or deploy intraday funds.
  \end{itemize}
\item \textbf{Lack of real-time data}, which complicates the task of determining and forecasting cash throughout the day.
\end{itemize}

% =====================================================================
\lo{LTR-6 b}{Discuss the governance structure of intraday liquidity risk management.}

Characteristics of an effective governance structure for overseeing intraday liquidity risk:

\begin{center}
\small
\begin{tabularx}{\textwidth}{@{}p{4.1cm} X@{}}
\toprule
\rowcolor{FRMSoft}\textbf{Characteristic} & \textbf{What it requires} \\
\midrule
\textbf{Active risk management} &
Intraday liquidity risk has too often been accepted as a cost of doing business and not actively managed. Effective firms instead classify \textbf{settlement and systemic risks} as components of their risk taxonomy and, critically, incorporate them into the firm's \textbf{risk appetite framework}. \\
\addlinespace
\textbf{Risk assessment} &
Treated as a component of risk self-assessments. \\
\addlinespace
\textbf{Integration with risk governance} &
Three lines of defence: \textbf{Treasury}, \textbf{risk management}, and \textbf{internal audit}. \\
\addlinespace
\textbf{Risk measurement and monitoring} &
The amount of intraday credit the institution is \textbf{extending to clients}, and the amount of intraday credit the institution \textbf{utilizes} itself. \\
\bottomrule
\end{tabularx}
\end{center}

\begin{lrtrapbox}[Who is which line of defence]
For intraday liquidity, the three lines are \textbf{Treasury} (first line, owns the risk and runs the position), \textbf{risk management} (second line, oversight), and \textbf{internal audit} (third line, independent assurance). Governance questions are won or lost on \emph{who}, not \emph{what} (\S5B.2, operator 3). Putting risk management in the first line, or naming the ALCO or the business lines instead of Treasury, is the standard build.
\end{lrtrapbox}

\begin{lrtrapbox}[Two directions of intraday credit, not one]
The measurement objective is symmetric: credit the bank \textbf{extends to clients} \emph{and} credit the bank \textbf{uses} itself. An option naming only one half is incomplete rather than wrong, which is exactly what a ``most appropriate'' qualifier is testing.
\end{lrtrapbox}

% =====================================================================
\lo{LTR-6 c}{Differentiate between methods for tracking intraday flows and monitoring risk levels.}

The objective is a \emph{differentiation}: two distinct families of measure, one counting flows and one quantifying risk. Keep them apart.

\begin{center}
\small
\begin{tabularx}{\textwidth}{@{}p{0.465\textwidth} p{0.465\textwidth}@{}}
\toprule
\rowcolor{PaleNavy}\textbf{Measures for \emph{tracking intraday flows}} &
\cellcolor{PaleGold}\textbf{Measures for \emph{quantifying and monitoring risk levels}} \\
\midrule
\textbf{Total payments} — maintain statistics concerning the amount of payments. &
\textbf{Daily maximum intraday liquidity usage.} \\
\addlinespace
\textbf{Other cash transactions} — track intraday and end-of-day settlement positions at all FMUs. &
\textbf{Intraday credit relative to Tier 1 capital.} \\
\addlinespace
\textbf{Time-sensitive obligations} — similar to settlement positions, these transactions require completion at a specific time during the day. &
\textbf{Client intraday credit usage} — derived by comparing a client's peak daily intraday overdraft to the established (committed or uncommitted) credit line. \\
\addlinespace
\textbf{Total intraday credit lines to clients and counterparties.} &
\textbf{Payment throughput} — tracks the percentage of outgoing payment activity relative to time of day. \\
\addlinespace
\textbf{Total bank intraday credit lines available and usage.} & \\
\bottomrule
\end{tabularx}
\end{center}
\figcap{Flow measures count \emph{what moved}; risk measures scale that against \emph{a limit or a capital base}. The tell is the denominator: every item in the right-hand column is a ratio or a peak-against-limit comparison, and every item in the left-hand column is a raw count or position.}

\begin{lrtrapbox}[Two pairs that are designed to be swapped]
\begin{itemize}
\item \textbf{Total intraday credit lines} (tracking, left) versus \textbf{client intraday credit usage} (risk, right). Both are about credit lines. The first counts the lines; the second compares a \emph{peak overdraft} to the line.
\item \textbf{Total payments} (tracking, left) versus \textbf{payment throughput} (risk, right). Both are about payments. The first is an amount; the second is a \emph{percentage relative to time of day}.
\end{itemize}
\end{lrtrapbox}

\begin{lrtrapbox}[Consolidated trap summary — LTR-6]
\begin{itemize}
\item \textbf{Polarity / role.} Largest \emph{use} $=$ outgoing wire payments. Largest \emph{source} $=$ incoming funds flow.
\item \textbf{Scope.} Incoming flows are the largest source \emph{in periods of normal market function}. The qualifier is the point.
\item \textbf{Scope.} Intraday credit lines are \textbf{uncommitted} and carry \textbf{no interest charge}.
\item \textbf{Role.} Three lines of defence here are Treasury, risk management, internal audit — in that order.
\item \textbf{Sibling.} Tracking measures are counts and positions; monitoring measures are ratios against a limit or Tier 1 capital. Total payments $\leftrightarrow$ payment throughput, and total credit lines $\leftrightarrow$ client credit usage, are the two swap pairs.
\item \textbf{Definition.} Client intraday credit usage compares the client's \textbf{peak} daily overdraft to the established line — not the average, and not the end-of-day balance.
\end{itemize}
\end{lrtrapbox}
